% !TEX root = chapter-standalone.tex

\section{Ordering monotone relations}
\label{sec:ordering-monotone-relations}

For a fixed source and target, the set of monotone functions can itself be made into a preorder in a straightforward manner (\cref{def:order-monotone-maps}): we order monotone functions ``point-wise''.
This construction extends to monotone relations, where it becomes even simpler: monotone relations are subsets, and we order them by inclusion.

\begin{remark}
    For \SY{preorders} \posA and \posB, the point-wise order on monotone functions $\posA \mto \posB$ is reflexive and transitive, because $\posBleq$ is; so it is a \SY{preorder}.
    It is a \SY{poset} when \posB is a poset, since then $\mapa(\posAel) \posBleq \mapb(\posAel)$ and $\mapb(\posAel) \posBleq \mapa(\posAel)$ for all $\posAel$ imply $\mapa = \mapb$.
\end{remark}

\begin{ctdefinition}[Order on monotone relations]
    \label{def:order-monotone-relations}
    Let \posA and \posB be \SY{preorders}, and write $\setOfMonRels(\posA, \posB)$ for the set of monotone relations $\posA \profto \posB$.
    We order $\setOfMonRels(\posA, \posB)$ by inclusion:
    \begin{equation}
        \prfdoubleperiod{
            \relA \posleqof{\setOfMonRels(\posA, \posB)} \relB
        }{
            \relA \setsubseteq \relB
        }
    \end{equation}
\end{ctdefinition}

Inclusion of subsets is reflexive, transitive, and antisymmetric.
So $\tup{\setOfMonRels(\posA, \posB), \setsubseteq}$ is always a \SY{poset}, even when \posA and \posB are only \SY{preorders}.
(Compare with the remark above: for monotone functions, we only get a poset when the target is a poset.)
We read $\relA \setsubseteq \relB$ as ``\relB relates everything that \relA relates, and possibly more''.

\begin{example}
    In \cref{exa:monrel-coins}, suppose that the shop also accepts payment by card for amounts up to five francs.
    The relation ``the amount $\ela$ can be paid with the coins $\tup{\elnc{1}, \elnc{2}}$ plus, possibly, a card payment of at most five francs'' is
    \begin{equation}
        \inrel{\ela}{\relB}{\tup{\elnc{1}, \elnc{2}}}
        \quad \text{if and only if} \quad
        \ela \leq \elnc{1} + 2\elnc{2} + 5.
    \end{equation}
    It is monotone for the same reasons as \relA, and $\relA \setsubseteq \relB$: every amount that can be paid with coins alone can also be paid with coins and a card.
\end{example}

\subsection{Unions and intersections}

The poset $\tup{\setOfMonRels(\posA, \posB), \setsubseteq}$ has a lot of structure: unions and intersections of monotone relations are again monotone relations.

\begin{lemma}
    \label{lem:monrel-unions-intersections}
    Let \posA and \posB be \SY{preorders}, and let $\tup{\relA_i}_{i \setin \setI}$ be a family of monotone relations $\posA \profto \posB$, indexed by some set \setI.
    Then the union $\bigsetunion_{i \setin \setI} \relA_i$ and the intersection $\bigsetintersection_{i \setin \setI} \relA_i$ (where, for $\setI = \Emptyset$, the intersection is taken to be $\posAset \cartprod \posBset$) are monotone relations $\posA \profto \posB$.
\end{lemma}
\begin{proof}
    For the union: suppose $\posAel' \posAleq \posAel$ and $\tup{\posAel, \posBel} \setin \bigsetunion_{i} \relA_i$.
    Then $\tup{\posAel, \posBel} \setin \relA_k$ for some $k \setin \setI$.
    Since $\relA_k$ is closed downward in the source, $\tup{\posAel', \posBel} \setin \relA_k \setsubseteq \bigsetunion_i \relA_i$.
    Closure upward in the target is proved in the same way.

    For the intersection: suppose $\posAel' \posAleq \posAel$ and $\tup{\posAel, \posBel} \setin \bigsetintersection_{i} \relA_i$.
    Then $\tup{\posAel, \posBel} \setin \relA_i$ for every $i \setin \setI$.
    Since each $\relA_i$ is closed downward in the source, $\tup{\posAel', \posBel} \setin \relA_i$ for every $i$, that is, $\tup{\posAel', \posBel} \setin \bigsetintersection_i \relA_i$.
    Closure upward in the target is proved in the same way.
    For $\setI = \Emptyset$, the intersection is the full relation, which is monotone.
\end{proof}

As a consequence, any collection of monotone relations $\posA \profto \posB$ has a least upper bound in $\tup{\setOfMonRels(\posA, \posB), \setsubseteq}$, namely its union, and a greatest lower bound, namely its intersection (compare with \cref{def:power-set-as-lattice}).
In particular, the empty relation is the least element and the full relation $\posAset \cartprod \posBset$ is the greatest element of $\setOfMonRels(\posA, \posB)$.

\subsection{Composition respects the order}

\begin{lemma}
    \label{lem:monrel-composition-monotone}
    Let $\relA, \relA' \colon \posA \profto \posB$ and $\relB, \relB' \colon \posB \profto \posC$ be monotone relations with $\relA \setsubseteq \relA'$ and $\relB \setsubseteq \relB'$.
    Then $\relA \mthen \relB \setsubseteq \relA' \mthen \relB'$.
\end{lemma}
\begin{proof}
    Suppose $\inrel{\posAel}{(\relA \mthen \relB)}{\posCel}$.
    Then there is $\posBel \setin \posBset$ with $\inrel{\posAel}{\relA}{\posBel}$ and $\inrel{\posBel}{\relB}{\posCel}$.
    Since $\relA \setsubseteq \relA'$ and $\relB \setsubseteq \relB'$, also $\inrel{\posAel}{\relA'}{\posBel}$ and $\inrel{\posBel}{\relB'}{\posCel}$.
    So the same $\posBel$ witnesses $\inrel{\posAel}{(\relA' \mthen \relB')}{\posCel}$.
\end{proof}

In other words, composition is a monotone function $\setOfMonRels(\posA, \posB) \Ptimes \setOfMonRels(\posB, \posC) \mto \setOfMonRels(\posA, \posC)$.
Similarly, the product and the sum of \cref{sec:monrel-new-from-old} are monotone in each argument, and the transpose is monotone as a function $\setOfMonRels(\posA, \posB) \mto \setOfMonRels(\posB\posop, \posAop)$; in each case this follows directly from the definitions.

\subsection{Companions, conjoints, and the point-wise order}

How does the order on monotone relations relate to the point-wise order on monotone functions?
The answer is that conjoints preserve the order, while companions reverse it.

\begin{lemma}
    \label{lem:companion-conjoint-order}
    Let $\mapa, \mapb \colon \posA \mto \posB$ be \SY{monotone functions} between \SY{preorders}.
    The following are equivalent:
    \begin{enumerate}
        \item $\mapa \posleqof{\posA \mto \posB} \mapb$, that is, $\mapa(\posAel) \posBleq \mapb(\posAel)$ for all $\posAel \setin \posAset$;
        \item $\comp{\mapb} \setsubseteq \comp{\mapa}$;
        \item $\conj{\mapa} \setsubseteq \conj{\mapb}$.
    \end{enumerate}
\end{lemma}
\begin{proof}
    \emph{(1)$\Imp$(2).}
    Suppose $\inrel{\posAel}{\comp{\mapb}}{\posBel}$, that is, $\mapb(\posAel) \posBleq \posBel$.
    Then $\mapa(\posAel) \posBleq \mapb(\posAel) \posBleq \posBel$, so $\inrel{\posAel}{\comp{\mapa}}{\posBel}$ by transitivity.

    \emph{(2)$\Imp$(1).}
    Let $\posAel \setin \posAset$.
    By reflexivity, $\inrel{\posAel}{\comp{\mapb}}{\mapb(\posAel)}$, so $\inrel{\posAel}{\comp{\mapa}}{\mapb(\posAel)}$, that is, $\mapa(\posAel) \posBleq \mapb(\posAel)$.

    \emph{(1)$\Imp$(3).}
    Suppose $\inrel{\posBel}{\conj{\mapa}}{\posAel}$, that is, $\posBel \posBleq \mapa(\posAel)$.
    Then $\posBel \posBleq \mapa(\posAel) \posBleq \mapb(\posAel)$, so $\inrel{\posBel}{\conj{\mapb}}{\posAel}$ by transitivity.

    \emph{(3)$\Imp$(1).}
    Let $\posAel \setin \posAset$.
    By reflexivity, $\inrel{\mapa(\posAel)}{\conj{\mapa}}{\posAel}$, so $\inrel{\mapa(\posAel)}{\conj{\mapb}}{\posAel}$, that is, $\mapa(\posAel) \posBleq \mapb(\posAel)$.
\end{proof}

The reversal for companions has a simple intuition: if $\mapb$ takes larger values than $\mapa$, then the condition $\mapb(\posAel) \posBleq \posBel$ is harder to satisfy than $\mapa(\posAel) \posBleq \posBel$, so $\comp{\mapb}$ relates fewer pairs.
Note also that \cref{lem:companion-conjoint-order} contains \cref{lem:companion-determines-function}: applying it in both directions shows that $\comp{\mapa} = \comp{\mapb}$ if and only if $\mapa \posleqof{\posA \mto \posB} \mapb$ and $\mapb \posleqof{\posA \mto \posB} \mapa$.

\begin{example}[Rounding functions]
    In \cref{ex:rounding-functions} we saw that $\funfloor \posleq \funceil$ as monotone functions $\realswithleq \mto \natswithleq$.
    By \cref{lem:companion-conjoint-order}, $\comp{\funceil} \setsubseteq \comp{\funfloor}$.
    Concretely, for $\ela \setin \reals$ and $\elb \setin \natnumbers$,
    \begin{equation}
        \inrel{\ela}{\comp{\funceil}}{\elb}
        \ \Eqv\
        \funceil(\ela) \leq \elb
        \ \Eqv\
        \ela \leq \elb,
    \end{equation}
    \begin{equation}
        \inrel{\ela}{\comp{\funfloor}}{\elb}
        \ \Eqv\
        \funfloor(\ela) \leq \elb
        \ \Eqv\
        \ela < \elb + 1.
    \end{equation}
    Indeed, every pair with $\ela \leq \elb$ also satisfies $\ela < \elb + 1$, but not conversely: for instance $\inrel{0.5}{\comp{\funfloor}}{0}$, while $0.5 \leq 0$ does not hold.
\end{example}

Finally, the order on monotone relations corresponds to the point-wise order on their curried forms (\cref{lem:monrel-currying}), once more with a reversal: for monotone relations $\relA, \relB \colon \posA \profto \posB$,
\begin{equation}
    \relA \setsubseteq \relB
    \quad \text{if and only if} \quad
    \hat{\relA}(\posAel) \setsubseteq \hat{\relB}(\posAel) \ \text{ for all } \posAel \setin \posAset,
\end{equation}
which says that $\hat{\relB} \posleq \hat{\relA}$ point-wise, as monotone functions into $\tup{\setOfUppersets \posB, \setsupseteq}$.
Indeed, $\tup{\posAel, \posBel} \setin \relA$ means $\posBel \setin \hat{\relA}(\posAel)$, and similarly for \relB.
